\documentclass[11pt]{article}
\usepackage[T1]{fontenc}
\usepackage{newtxtext,newtxmath}
\usepackage[a4paper,top=2.0cm,bottom=2.0cm,left=2.4cm,right=2.4cm]{geometry}
\usepackage{enumitem}
\usepackage{url}
\usepackage{xcolor}
% Author details, shared by main.tex, si.tex and cover_letter.tex. Edit only this file.
\newcommand{\AuthorName}{Albert Linda}
\newcommand{\AuthorAffiliation}{Department of Materials Science and Engineering, Indian Institute of Technology Kanpur, Kalyanpur, Kanpur 208016, Uttar Pradesh, India}
\newcommand{\AuthorORCID}{0000-0002-3723-1756}
\newcommand{\AuthorEmail}{albert.hzbn@gmail.com}
\newcommand{\Funding}{This research received no specific grant from any funding agency in the public, commercial or not-for-profit sectors.}
\newcommand{\AuthorAddressLines}{Department of Materials Science and Engineering\\Indian Institute of Technology Kanpur\\Kalyanpur, Kanpur 208016, Uttar Pradesh, India}

\pagestyle{empty}
\setlength{\parindent}{0pt}
\setlength{\parskip}{0.55\baselineskip}
\urlstyle{same}

\begin{document}

% ------------------------------------------------------------------ letterhead
{\large\bfseries \AuthorName}\\[2pt]
Department of Materials Science and Engineering\\
Indian Institute of Technology Kanpur, Kanpur 208016, Uttar Pradesh, India\\
\AuthorEmail\ \ $\cdot$\ \ ORCID \AuthorORCID

\vspace{-4pt}
{\color{black!60}\rule{\textwidth}{0.4pt}}

\begin{minipage}[t]{0.6\textwidth}
Editor-in-Chief\\
\emph{Journal of Chemical Theory and Computation}\\
American Chemical Society
\end{minipage}\hfill
\begin{minipage}[t]{0.35\textwidth}\raggedleft
\today
\end{minipage}

\textbf{Re: Submission of the manuscript ``Factorised Cartesian Atomic Cluster Expansion for Low-Cost
Machine-Learning Interatomic Potentials''}

Dear Editor,

Please consider the enclosed manuscript for publication as an Article in the \emph{Journal of Chemical Theory and
Computation}.

Machine-learned interatomic potentials of the atomic cluster expansion (ACE) family, such as MACE, reach near-DFT
accuracy, but their cost per evaluation still limits molecular-dynamics length and system size. The manuscript makes
four contributions within the ACE framework. It does not claim a new theoretical foundation.
\begin{enumerate}[leftmargin=1.6em, itemsep=2pt, topsep=0pt, parsep=0pt]
\item \textbf{Factorised edge density.} A factorised Cartesian form of the ACE edge density whose per-edge cost
does not depend on the feature width. Tensor transport is a single constant linear map, and the body-ordered algebra
runs in fused GPU kernels. At equal training budget the model is more accurate than MACE-small on four of five rMD17
molecules and tied on the fifth, and it is 5.7 times faster at inference than MACE-small and 10.9 times faster than
MACE-medium.
\item \textbf{Physics operators.} One interface through which latent-charge electrostatics (Ewald and smooth
particle-mesh Ewald), dispersion and short-range repulsion are attached, adding about 23\% to the cost of a 4,500-atom system. It is validated against
analytical and Ewald references and on the charged-system benchmarks of Ko et al.
\item \textbf{Training stability.} A diagnosis and fix for the training divergences of this model class (unit
weight parameterisation), supported by a controlled probe.
\item \textbf{Distillation.} Uncertainty-gated, iterated ensemble distillation. On rMD17 aspirin it brings a single
model to energy and force errors of 2.26~meV and 6.74~meV/\AA, within 3\% and 2\% of the published MACE result, at
lower inference cost than a 64-channel MACE.
\end{enumerate}

We report the limitations openly. Training steps are no cheaper than for MACE-small, and distillation multiplies
training compute. Several recent equivariant models remain more accurate on aspirin. All MACE comparisons use the
reference implementation without the cuEquivariance kernels.

The work suits JCTC's readership in method development for molecular simulation. All code, tests, training logs,
checkpoints and the scripts that regenerate every table and figure are publicly available under the MIT licence at
\url{https://github.com/albert-hzbn/face-mlip}. An AI assistant was used during the work, as disclosed in the
manuscript.

This manuscript is not under consideration elsewhere and has not been published previously. The author declares no
competing financial interest.

Sincerely,

\vspace{1.2cm}
\AuthorName\\
Department of Materials Science and Engineering\\
Indian Institute of Technology Kanpur

\end{document}
